\documentclass[../main.tex]{subfiles}
\ifSubfilesClassLoaded{\setcounter{chapter}{5}}{}
\begin{document}

\chapter{EDA Univariat dan Bivariat; Matplotlib dan Seaborn}

\begin{subcpmk}
  \item \textbf{Sub-CPMK-P4:} EDA dan visualisasi pola dengan minimal empat jenis plot berbeda serta interpretasi singkat (RPS Mg.\ 5).
\end{subcpmk}

\SesiRPS{5}{P4}{$\sim$170 menit}{CPMK-P2}

\section{Kemampuan akhir sesi}

Menghasilkan histogram/KDE, scatter atau boxplot kategori, heatmap korelasi; setiap plot memiliki judul dan label sumbu; merumuskan satu hipotesis pemodelan.

\section{Teori}

\begin{konsep}
EDA menjawab ``apa yang data ceritakan'' sebelum asumsi model. Visual yang baik memiliki \textbf{judul}, \textbf{label sumbu}, dan \textbf{skala} yang terbaca; hindari chart junk berlebihan.
\end{konsep}

\begin{catatan}
Contoh memakai dataset \texttt{titanic} Seaborn; target biner \texttt{survived} cocok untuk hipotesis klasifikasi (mis.\ hubungan \texttt{fare}/\texttt{class} dengan survival).
\end{catatan}

\section{Sarana}

\texttt{matplotlib.pyplot}, \texttt{seaborn}, \texttt{pandas}.

\section{Contoh kode}

\begin{lstlisting}[caption=Histogram umur dan scatter umur vs fare]
import matplotlib.pyplot as plt
import seaborn as sns

df = sns.load_dataset("titanic")

fig, ax = plt.subplots(1, 2, figsize=(10, 4))
sns.histplot(df["age"].dropna(), kde=True, ax=ax[0])
ax[0].set_title("Distribusi umur penumpang")
ax[0].set_xlabel("age")
sns.scatterplot(
    data=df, x="age", y="fare", hue="survived", alpha=0.65, ax=ax[1]
)
ax[1].set_title("Umur vs tarif (hue: survived)")
plt.tight_layout()
plt.show()
\end{lstlisting}

\begin{lstlisting}[caption=Boxplot tarif per kelas kabin]
plt.figure(figsize=(8, 4))
sns.boxplot(data=df, x="class", y="fare")
plt.title("Sebaran fare per class")
plt.xlabel("class")
plt.ylabel("fare")
plt.tight_layout()
plt.show()
\end{lstlisting}

\begin{lstlisting}[caption=Heatmap korelasi fitur numerik]
num = df[["survived", "pclass", "age", "sibsp", "parch", "fare"]].dropna()
plt.figure(figsize=(7, 5))
sns.heatmap(num.corr(), annot=True, fmt=".2f", cmap="vlag", center=0)
plt.title("Korelasi antar variabel numerik")
plt.tight_layout()
plt.show()
\end{lstlisting}

\begin{lstlisting}[caption=Pairplot subset (dengan hue survival)]
cols = ["age", "fare", "pclass"]
sns.pairplot(
    df[cols + ["survived"]].dropna(),
    hue="survived",
    corner=True,
    plot_kws={"alpha": 0.6},
)
plt.show()
\end{lstlisting}

\section{Referensi}

\begin{itemize}[leftmargin=*]
  \item Microsoft Explore data. \url{https://learn.microsoft.com/en-us/training/modules/explore-analyze-data-with-python/}
  \item Seaborn tutorial. \url{https://seaborn.pydata.org/tutorial.html}
\end{itemize}

\begin{asesmen}
\textbf{Tugas modul:} empat jenis visual berbeda + satu kalimat interpretasi per plot + satu hipotesis model.
\end{asesmen}

\begin{rangkuman}
Minggu 5 menghubungkan statistik deskriptif dengan visual untuk pertanyaan analitik.

\textbf{Kata kunci:} histogram, boxplot, heatmap, pairplot
\end{rangkuman}

\end{document}
